% Part: first-order-logic
% Chapter: syntax-and-semantics
% Section: main-operator

\documentclass[../../../include/open-logic-section]{subfiles}

\begin{document}

\olfileid{fol}{syn}{mai}

\olsection{ఒక సూత్రపు \printtoken{S}{main operator}}

\begin{explain}
\tetoken{ఒక సూత్రం}{!!a{formula}}~$!A$ను నిర్మించడంలో చివరిగా ఉపయోగించిన సంయోజకం గురించి
మాట్లాడడం తరచుగా ఉపయోగకరం. ఈ సంయోజకాన్ని $!A$ యొక్క
\emph{ప్రధాన సంయోజకం} అంటాం. సహజావబోధన ప్రకారం, ఇది $!A$లో
``అత్యంత వెలుపలి'' సంయోజకం. ఉదాహరణకు, $\lnot !A$కు ప్రధాన సంయోజకం
$\lnot$; $(!A \lor !B)$కు ప్రధాన సంయోజకం $\lor$; ఇలాగే కొనసాగుతుంది.
\end{explain}


\begin{defn}[\tetoken{ప్రధాన సంయోజకం}{!!^{main operator}}]
\ollabel{def:main-op}
\tetoken{ఒక సూత్రం}{!!a{formula}}~$!A$ యొక్క \emph{\tetoken{ప్రధాన సంయోజకం}{!!{main operator}}}ను కింది విధంగా
నిర్వచిస్తాం:
\begin{enumerate}
\item \indcase*{!A}{!A}{$\indfrm$కు \tetoken{ప్రధాన సంయోజకం}{!!{main operator}} ఉండదు.}

\tagitem{prvNot}{\indcase{!A}{\lnot !B}{$\indfrm$ యొక్క \tetoken{ప్రధాన సంయోజకం}{!!{main operator}}
  $\lnot$.}}{}

\tagitem{prvAnd}{\indcase{!A}{(!B \land !C)}{$\indfrm$ యొక్క
  \tetoken{ప్రధాన సంయోజకం}{!!{main operator}} $\land$.}}{}

\tagitem{prvOr}{\indcase{!A}{(!B \lor !C)}{$\indfrm$ యొక్క
  \tetoken{ప్రధాన సంయోజకం}{!!{main operator}} $\lor$.}}{}

\tagitem{prvIf}{\indcase{!A}{(!B \lif !C)}{$\indfrm$ యొక్క
  \tetoken{ప్రధాన సంయోజకం}{!!{main operator}} $\lif$.}}{}

\tagitem{prvIff}{\indcase{!A}{(!B \liff !C)}{$\indfrm$ యొక్క
  \tetoken{ప్రధాన సంయోజకం}{!!{main operator}} $\liff$.}}{}

\tagitem{prvAll}{\indcase{!A}{\lforall[x][!B]}{$\indfrm$ యొక్క
  \tetoken{ప్రధాన సంయోజకం}{!!{main operator}} $\lforall$.}}{}

\tagitem{prvEx}{\indcase{!A}{\lexists[x][!B]}{$\indfrm$ యొక్క
  \tetoken{ప్రధాన సంయోజకం}{!!{main operator}} $\lexists$.}}{}
\end{enumerate}
\end{defn}

ప్రతి సందర్భంలోనూ సూత్రంలో సూచించిన నిర్దిష్ట \tetoken{ప్రధాన సంయోజకం}{!!{main operator}}
\emph{ఘటన}నే ఉద్దేశిస్తున్నాం. ఉదాహరణకు,
$((!D \lif !E) \lif (!E \lif !D))$ సూత్రం $(!B \lif !C)$ రూపంలో
ఉంటుంది; ఇక్కడ $!B$ అనేది $(!D \lif !E)$, $!C$ అనేది
$(!E \lif !D)$. కనుక $\lif$ యొక్క రెండవ ఘటనే \tetoken{ప్రధాన సంయోజకం}{!!{main operator}}.

\begin{explain}
ఇది పరమాణుయేతర \tetoken{సూత్రాలు}{!!{formula}s} అన్నింటినీ వాటి \tetoken{ప్రధాన సంయోజకం}{!!{main operator}} ఘటనకు
పంపే ప్రమేయానికి \emph{పునరావర్తక} నిర్వచనం. \tetoken{సూత్రాలను}{!!{formula}s}
ఆగమనాత్మకంగా నిర్వచించిన తీరు వల్ల ప్రతి \tetoken{సూత్రం}{!!{formula}}~$!A$,
\olref{def:main-op}లోని సందర్భాల్లో ఒకదాన్ని తీరుస్తుంది. అందువల్ల
ప్రతి పరమాణుయేతర \tetoken{సూత్రం}{!!{formula}}~$!A$కు \tetoken{ఒక ప్రధాన సంయోజకం}{!!a{main
  operator}} ఉంటుందని హామీ.
ప్రతి \tetoken{సూత్రం}{!!{formula}} ఈ షరతుల్లో ఒక్కదాన్నే తీరుస్తుంది; ప్రతి సందర్భంలో
$!A$ను నిర్మించిన చిన్న \tetoken{సూత్రాలు}{!!{formula}s} ఏకైకంగా నిర్ణయితమవుతాయి. కనుక
$!A$ యొక్క \tetoken{ప్రధాన సంయోజకం}{!!{main
  operator}} ఘటన ఏకైకం; అందువల్ల మనం ఒక ప్రమేయాన్ని
నిర్వచించాం.
\end{explain}

ఒక సూత్రపు \tetoken{ప్రధాన సంయోజకం}{!!{main operator}} ఏ సంకేతమో బట్టి, \olref{tab:main-op}లోని
పేర్లతో \tetoken{సూత్రాలను}{!!{formula}s} పిలుస్తాం.\iftag{defNot,defOr,defAnd,defIf,defIff,defTrue,defFalse,defEx,defAll}
{ అయితే నిర్వచిత సంయోజకాలు అధికారికంగా \tetoken{సూత్రాలలో}{!!{formula}s} కనిపించవని
గుర్తుంచుకోండి. అవి కేవలం సంక్షిప్త రూపాలు; అందువల్ల అధికారికంగా అవి
ఒక సూత్రానికి ప్రధాన సంయోజకాలు కాలేవు. కాబట్టి అన్ని \tetoken{సూత్రాలు}{!!{formula}s}
గురించిన నిరూపణల్లో వాటిని విడిగా పరిశీలించాల్సిన అవసరం లేదు.}

\begin{table}[!h]
\centering
\begin{tabular}{c | c | c}
\tetoken{ప్రధాన సంయోజకం}{!!^{main operator}} & \tetoken{సూత్రం}{!!{formula}} రకం & ఉదాహరణ\\
\hline
ఏదీ లేదు & పరమాణు (\tetoken{సూత్రం}{!!{formula}}) &
\iftag{prvFalse}{$\lfalse$,}{}
\iftag{prvTrue}{$\ltrue$,}{}
$\Atom{R}{t_1, \dots, t_n}$,
$\eq[t_1][t_2]$\\
$\lnot$ & నిషేధం & $\lnot !A$ \\
$\land$ & సంధానం & $(!A \land !B)$ \\
$\lor$ & వికల్పం & $(!A \lor !B)$ \\
$\lif$ & \tetoken{సోపాధికం}{!!{conditional}} & $(!A \lif !B)$ \\
$\liff$ & \tetoken{ద్విసోపాధికం}{!!{biconditional}} & $(!A \liff !B)$ \\
$\lforall[][]$ & సార్వత్రిక (\tetoken{సూత్రం}{!!{formula}})& $\lforall[x][!A]$ \\
$\lexists[][]$ & అస్తిత్వ (\tetoken{సూత్రం}{!!{formula}})& $\lexists[x][!A]$
\end{tabular}
\caption{\tetoken{సూత్రాలు}{!!{formula}s} ప్రధాన సంయోజకం, పేర్లు}
\ollabel{tab:main-op}
\end{table}
\sourcecorrection{OLTEFOLSYN-002}{ఆంగ్ల మూలంలో సంధాన ఉదాహరణకు కుడి
కుండలీకరణ చిహ్నం గణిత పరిధి వెలుపల ఉంది. దాన్ని $(!A \land !B)$లో
జతగా ఉండే స్థానానికి మార్చాం; సూత్రంలోని సంకేతాల క్రమం మారలేదు.}

\end{document}
